
% LaTeX Beamer file automatically generated from DocOnce
% https://github.com/doconce/doconce

%-------------------- begin beamer-specific preamble ----------------------

\documentclass{beamer}

\usetheme{red_plain}
\usecolortheme{default}

% turn off the almost invisible, yet disturbing, navigation symbols:
\setbeamertemplate{navigation symbols}{}

% Examples on customization:
%\usecolortheme[named=RawSienna]{structure}
%\usetheme[height=7mm]{Rochester}
%\setbeamerfont{frametitle}{family=\rmfamily,shape=\itshape}
%\setbeamertemplate{items}[ball]
%\setbeamertemplate{blocks}[rounded][shadow=true]
%\useoutertheme{infolines}
%
%\usefonttheme{}
%\useinntertheme{}
%
%\setbeameroption{show notes}
%\setbeameroption{show notes on second screen=right}

% fine for B/W printing:
%\usecolortheme{seahorse}

\usepackage{pgf}
\usepackage{graphicx}
\usepackage{epsfig}
\usepackage{relsize}

\usepackage{fancybox}  % make sure fancybox is loaded before fancyvrb

\usepackage{fancyvrb}
\usepackage{minted} % requires pygments and latex -shell-escape filename
%\usepackage{anslistings}
%\usepackage{listingsutf8}

\usepackage{amsmath,amssymb,bm}
%\usepackage[latin1]{inputenc}
\usepackage[T1]{fontenc}
\usepackage[utf8]{inputenc}
\usepackage{colortbl}
\usepackage[english]{babel}
\usepackage{tikz}
\usepackage{framed}
% Use some nice templates
\beamertemplatetransparentcovereddynamic

% --- begin table of contents based on sections ---
% Delete this, if you do not want the table of contents to pop up at
% the beginning of each section:
% (Only section headings can enter the table of contents in Beamer
% slides generated from DocOnce source, while subsections are used
% for the title in ordinary slides.)
\AtBeginSection[]
{
  \begin{frame}<beamer>[plain]
  \frametitle{}
  %\frametitle{Outline}
  \tableofcontents[currentsection]
  \end{frame}
}
% --- end table of contents based on sections ---

% If you wish to uncover everything in a step-wise fashion, uncomment
% the following command:

%\beamerdefaultoverlayspecification{<+->}

\newcommand{\shortinlinecomment}[3]{\note{\textbf{#1}: #2}}
\newcommand{\longinlinecomment}[3]{\shortinlinecomment{#1}{#2}{#3}}

\definecolor{linkcolor}{rgb}{0,0,0.4}
\hypersetup{
    colorlinks=true,
    linkcolor=linkcolor,
    urlcolor=linkcolor,
    pdfmenubar=true,
    pdftoolbar=true,
    bookmarksdepth=3
    }
\setlength{\parskip}{0pt}  % {1em}

\newenvironment{doconceexercise}{}{}
\newcounter{doconceexercisecounter}
\newenvironment{doconce:movie}{}{}
\newcounter{doconce:movie:counter}

\newcommand{\subex}[1]{\noindent\textbf{#1}}  % for subexercises: a), b), etc

\logo{{\tiny \copyright\ 1999-2025, Morten Hjorth-Jensen. Released under CC Attribution-NonCommercial 4.0 license}}

%-------------------- end beamer-specific preamble ----------------------

% Add user's preamble




% insert custom LaTeX commands...

\raggedbottom
\makeindex

%-------------------- end preamble ----------------------

\begin{document}

% matching end for #ifdef PREAMBLE

\newcommand{\exercisesection}[1]{\subsection*{#1}}



% ------------------- main content ----------------------



% ----------------- title -------------------------

\title{Quantum Computing, Quantum Machine Learning and Quantum Information Theories}

% ----------------- author(s) -------------------------

\author{Morten Hjorth-Jensen\inst{1}}
\institute{Department of Physics, University of Oslo, Norway\inst{1}}
% ----------------- end author(s) -------------------------

\date{February 19, 2025
% <optional titlepage figure>
\ \\ 
{\tiny \copyright\ 1999-2025, Morten Hjorth-Jensen. Released under CC Attribution-NonCommercial 4.0 license}
}

\begin{frame}[plain,fragile]
\titlepage
\end{frame}

\begin{frame}[plain,fragile]
\frametitle{Plans for the week of February 17-21, 2025}

\begin{block}{}
The first part of the lecture will focus
\begin{enumerate}
\item Review from last week, one-qubit gates and one-qubit Hamiltonian (using slides from last week)

\item Reminder and review of  density matrices and measurements from week 2

\item Schmidt decomposition and entanglement

\item Discussion of entropies, classical information entropy (Shannon entropy) and von Neumann entropy
\end{enumerate}

\noindent
The second part focuses on 
\begin{enumerate}
\item Two-qubit Hamiltonians and how to encode them

\item Discussion of project 1 during exercise session
% o \href{{https://youtu.be/}}{Video of lecture to be added}
% o \href{{https://github.com/CompPhysics/QuantumComputingMachineLearning/blob/gh-pages/doc/HandWrittenNotes/2024/NotesFebruary14.pdf}}{Whiteboard notes}
\end{enumerate}

\noindent
\end{block}
\end{frame}

\begin{frame}[plain,fragile]
\frametitle{Motivation}

In order to study entanglement and why it is so important for quantum
computing, we need to introduce some basic measures and useful
quantities.  For these endeavors, we  will use our two-qubit system from
the second lecture in order to introduce and repeat, through examples, density
matrices and entropy. These two quantities, together with
technicalities like the Schmidt decomposition define important quantities in analyzing quantum computing examples.

The Schmidt decomposition is again a
linear decomposition which allows us to express a vector in terms of
tensor product of two inner product spaces. In quantum information
theory and quantum computing it is widely used as away to define and
describe entanglement.
\end{frame}

\begin{frame}[plain,fragile]
\frametitle{Reminder on density matrices and traces}

We have the spectral decomposition of a given operator $\bm{A}$ given by

\[
\bm{A}=\sum_i\lambda_i\vert i \rangle\langle i\vert,
\]
with the ONB $\vert i\rangle$ being eigenvectors of $\bm{A}$ and $\lambda_i$ being the eigenvalues. Similarly, a operator which is a function of $\bm{A}$ is given by
\[
f(\bm{A})=\sum_if(\bm{A})\vert i \rangle\langle i\vert.
\]
The trace of a product of matrices is cyclic, that is
\[
\mathrm{tr}[\bm{ABC}])=\mathrm{tr}[\bm{BCA}])=\mathrm{tr}[\bm{CBA}]),
\]
and we have also 
\[
\mathrm{tr}[\bm{A}\vert \psi\rangle\langle\psi\vert])=\langle\psi\vert\bm{A}\vert\psi\rangle.
\]
\end{frame}

\begin{frame}[plain,fragile]
\frametitle{Definition of density matrix}

Using the spectral decomposition we defined also the density matrix as
\[
\rho = \sum_i p_i\vert i \rangle\langle i\vert,
\]
where the probability $p_i$ are the eigenvalues of the density linked with the ONB $\vert i \rangle$.

The trace of the density matrix is $\mathrm{tr}\rho=1$ and it is invariant under unitary transformations
$\vert \psi_i'\rangle = \bm{U}\vert \psi_i\rangle$.
The unitary transformation of the density matrix gives, with
$\bm{U}^{\dagger}\bm{U}=\bm{U}^{T}\bm{U}=\bm{I}$,
\[
\bm{U}\rho\bm{U}^{\dagger}=\sum_ipi_i\bm{U}\vert \psi_i\rangle\langle\psi_i\vert\bm{U}^{\dagger},
\]
and with the unitary transformation it is easy to show that thet trace of the transformaed density matrix is equal to one,
\[
\mathrm{tr}\left[ \bm{U}\rho\bm{U}^{\dagger}\right]=\mathrm{tr}\left[ \bm{U}\bm{U}^{\dagger}\rho\right]=1.
\]
\end{frame}

\begin{frame}[plain,fragile]
\frametitle{Revisiting our Bell states and ur  entanglement encounter, two qubit system}

We define a system that can be thought of as composed of two subsystems
$A$ and $B$. Each subsystem has computational basis states

\[
\vert 0\rangle_{\mathrm{A,B}}=\begin{bmatrix} 1 & 0\end{bmatrix}^T \hspace{1cm} \vert 1\rangle_{\mathrm{A,B}}=\begin{bmatrix} 0 & 1\end{bmatrix}^T.
\]
The subsystems could represent single particles or composite many-particle systems of a given symmetry.
\end{frame}

\begin{frame}[plain,fragile]
\frametitle{Computational basis}

This leads to the many-body computational basis states

\[
\vert 00\rangle = \vert 0\rangle_{\mathrm{A}}\otimes \vert 0\rangle_{\mathrm{B}}=\begin{bmatrix} 1 & 0 & 0 &0\end{bmatrix}^T,
\]
and
\[
\vert 01\rangle = \vert 0\rangle_{\mathrm{A}}\otimes \vert 1\rangle_{\mathrm{B}}=\begin{bmatrix} 0 & 1 & 0 &0\end{bmatrix}^T,
\]
and
\[
\vert 10\rangle = \vert 1\rangle_{\mathrm{A}}\otimes \vert 0\rangle_{\mathrm{B}}=\begin{bmatrix} 0 & 0 & 1 &0\end{bmatrix}^T,
\]
and finally
\[
\vert 11\rangle = \vert 1\rangle_{\mathrm{A}}\otimes \vert 1\rangle_{\mathrm{B}}=\begin{bmatrix} 0 & 0 & 0 &1\end{bmatrix}^T.
\]
\end{frame}

\begin{frame}[plain,fragile]
\frametitle{Bell states}

The above computational basis states, which define an ONB, can in turn
be used to define another ONB. As an example, consider the so-called
Bell states

\[
\vert \Phi^+\rangle = \frac{1}{\sqrt{2}}\left[\vert 00\rangle +\vert 11\rangle\right]=\frac{1}{\sqrt{2}}\begin{bmatrix} 1 \\ 0 \\ 0 \\ 1\end{bmatrix},
\]

\[
\vert \Phi^-\rangle = \frac{1}{\sqrt{2}}\left[\vert 00\rangle -\vert 11\rangle\right]=\frac{1}{\sqrt{2}}\begin{bmatrix} 1 \\ 0 \\ 0 \\ -1\end{bmatrix},
\]
\end{frame}

\begin{frame}[plain,fragile]
\frametitle{The next two}

\[
\vert \Psi^+\rangle = \frac{1}{\sqrt{2}}\left[\vert 10\rangle +\vert 01\rangle\right]=\frac{1}{\sqrt{2}}\begin{bmatrix} 0 \\ 1 \\ 1 \\ 0\end{bmatrix},
\]

and

\[
\vert \Psi^-\rangle = \frac{1}{\sqrt{2}}\left[\vert 10\rangle -\vert 01\rangle\right]=\frac{1}{\sqrt{2}}\begin{bmatrix} 0 \\ 1 \\ -1 \\ 0\end{bmatrix}.
\]
It is easy to convince oneself that these states also form an orthonormal basis.
\end{frame}

\begin{frame}[plain,fragile]
\frametitle{Measurement}

Measuring one of the qubits of one of the above Bell states,
automatically determines, as we will see below, the state of the
second qubit. To convince ourselves about this, let us assume we perform a measurement on the qubit in system $A$ by introducing the projections with outcomes $0$ or $1$ as

\[
\bm{P}_0=\vert 0\rangle\langle 0\vert_A\otimes \bm{I}_B=\begin{bmatrix} 1 & 0\\ 0 & 0\end{bmatrix}\otimes\begin{bmatrix} 1& 0 \\ 0 & 1\end{bmatrix}=\begin{bmatrix} 1 & 0 & 0 & 0 \\ 0 & 1 & 0 & 0 \\ 0 & 0 & 0 & 0 \\ 0 & 0 & 0 & 0\end{bmatrix},
\]
for the projection of the $\vert 0 \rangle$ state in system $A$ and similarly
\[
\bm{P}_1=\vert 1\rangle\langle 1\vert_A\otimes \bm{I}_B=\begin{bmatrix} 0 & 0\\ 0 & 1\end{bmatrix}\otimes\begin{bmatrix} 1& 0 \\ 0 & 1\end{bmatrix}=\begin{bmatrix} 0 & 0 & 0 & 0 \\ 0 & 0 & 0 & 0 \\ 0 & 0 & 1 & 0 \\ 0 & 0 & 0 & 1\end{bmatrix},
\]
for the projection of the $\vert 1 \rangle$ state in system $A$.
\end{frame}

\begin{frame}[plain,fragile]
\frametitle{Probability of  outcome}

We can then calculate the probability for the various outcomes by
computing for example the probability for measuring qubit $0$ 

\[
\langle \Phi^+\vert \bm{P}_0\vert \Phi^+\rangle = \frac{1}{2} \left[\langle 00\vert +\langle 11\vert\right]\begin{bmatrix} 1 & 0 & 0 & 0 \\ 0 & 1 & 0 & 0 \\ 0 & 0 & 0 & 0 \\ 0 & 0 & 0 & 0\end{bmatrix}\left[\vert 00\rangle +\vert 11\rangle\right]=\frac{1}{2}.
\]
Similarly, we obtain
\[
\langle \Phi^+\vert \bm{P}_1\vert \Phi^+\rangle = \frac{1}{2}\left[\langle 00\vert +\langle 11\vert\right]\begin{bmatrix} 0 & 0 & 0 & 0 \\ 0 & 0 & 0 & 0 \\ 0 & 0 & 1 & 0 \\ 0 & 0 & 0 & 1\end{bmatrix}\left[\vert 00\rangle +\vert 11\rangle\right]=\frac{1}{2}.
\]
\end{frame}

\begin{frame}[plain,fragile]
\frametitle{States after measurement}

After the above measurements the system is in the states

\[
\vert \Phi'_0 \rangle = \sqrt{2}\left[\vert 0\rangle\langle 0\vert_A\otimes \bm{I}_B\right]\vert\Phi^+\rangle=\vert 00\rangle,
\]
and 
\[
\vert \Phi'_1 \rangle = \sqrt{2}\left[\vert 1\rangle\langle 1\vert_A\otimes \bm{I}_B\right]\vert\Phi^+\rangle=\vert 11\rangle.
\]

We see from the last two equations that the state of the second qubit
is determined even though the measurement has only taken place locally
on system $A$.
\end{frame}

\begin{frame}[plain,fragile]
\frametitle{Other states}

If we on the other hand consider a state like

\[
\vert 00 \rangle = \vert 0\rangle_A\otimes \vert 0\rangle_B,
\]
this is a pure \textbf{product} state of the single-qubit, or single-particle
states, of two qubits (particles) in system $A$ and system $B$,
respectively. We call such a state for a \textbf{pure product state}.  Quantum states
that cannot be written as a mixture of other states are called pure
quantum states or just product states, while all other states are called mixed quantum states.
\end{frame}

\begin{frame}[plain,fragile]
\frametitle{More on Bell states}

A state like one of the Bell states (where we introduce the subscript $AB$ to indicate that the state is composed of single states from two subsystem)
\[
\vert \Phi^+\rangle = \frac{1}{\sqrt{2}}\left[\vert 00\rangle_{AB} +\vert 11\rangle_{AB}\right],
\]
is on the other hand a mixed state and we cannot determine whether system $A$ is in a state $0$ or $1$. The above state is a superposition of the states $\vert 00\rangle_{AB}$ and $\vert 11\rangle_{AB}$ and it is not possible to determine individual states of systems $A$ and $B$, respectively.
\end{frame}

\begin{frame}[plain,fragile]
\frametitle{Entanglement}

We say that the state is entangled. This yields the following
definition of entangled states: a pure bipartite state $\vert
\psi\rangle_{AB}$ is entangled if it cannot be written as a product
state $\vert\psi\rangle_{A}\otimes\vert\phi\rangle_B$ for any choice
of the states $\vert\psi\rangle_{A}$ and $\vert\phi\rangle_B$. Otherwise we say the state is separable.
\end{frame}

\begin{frame}[plain,fragile]
\frametitle{Maximally entangled}

A so-called maximally entangled state for a bipartite system has equal  probability amplitudes
\[
\vert \Psi \rangle = \frac{1}{\sqrt{d}}\sum_{i=0}^{d-1}\vert ii\rangle.
\]

We call a bipartite state composed of systems $A$ and $B$ (these
systems can be single-particle systems, or single-qubit systems
representing low-lying states of complicated many-body systems) for
separable if its density matrix $\rho_{AB}$ can be written out as the
tensor product of the individual density matrices $\rho_A$ and
$\rho_B$, that is we have for a given probability distribution $p_i$

\[
\rho_{AB}=\sum_ip_i\rho_A(i)\otimes \rho_B(i).
\]
\end{frame}

\begin{frame}[plain,fragile]
\frametitle{Entropies and density matrices}

We discuss first the classical information entropy.
Thereafter we define the quantum-mechanical quantity, commonly known as the Von-Neumann entropy.
\end{frame}

\begin{frame}[plain,fragile]
\frametitle{Shannon information entropy}

We start our discussions with the classical information entropy, or
just Shannon entropy, before we move over to a quantum mechanical way
to define the entropy based on the density matrices discussed earlier.

We define a set of random variables $X=\{x_0,x_1,\dots,x_{n-1}\}$ with probability for an outcome $x\in X$ given by $p_X(x)$, the
classical information entropy is defined as
\[
S=-\sum_{x\in X}p_X(x)\log_2{p_X(x)}.
\]
Why this expression? What does it mean?
\end{frame}

\begin{frame}[plain,fragile]
\frametitle{Mathematics of entropy}

What is the basic idea of the entropy as it is used in information
theory (we leave out the standard description from statistical physics
here)?

We want to have a measure of unlikelihood. Consider a simple binary
system, with two outcomes, true and false with true given by a
probability $p$ and false given by $1-p$. Since $p$ represents a
probability, and assuming the probabilities are properly normalized,
we have $p\in [0,1]$.

Suppose know that we know precisely $p$, for example we could set
$p=1$. This means that there is only one outcome, namely the true one
and false is equal to zero. We know thus (exactly) which state the
system is in.  If we set $p=0$, then we know that the outcome is
false, without doubt.  Can we find a way to, through a specific
function to say with a given certainty that the system is in a
specific state? Let us call this function for $S$.
\end{frame}

\begin{frame}[plain,fragile]
\frametitle{More on the mathematics of entropy}

We could for example use
\[
S(x)= -\log_2{p(x)},
\]
as potential
function. Note that if we change the log-base, we only change the results by a scaling constant.

When $p(x)=1$, then $S(x)=0$ and we could say that the
probability of being in one specific state is uniquely defined since
$S$ is then zero.  This function is also continuous, it is additive
(which is very useful of we have independent and identically
distributed stochastic variables), it is also high for unlikely events
since when $p$ goes to zero (unlikely event) it becomes very
large. It is also a non-negative function.
\end{frame}

\begin{frame}[plain,fragile]
\frametitle{Changing the expression}

We note however that if we go back to our binary model and would like to
describe the entropy in terms of the variable $p$, we end up with $p$ being zero for the false outcome.
The latter has probability one ($p=0)$.
Using that
\[
\lim_{x\to -\infty} -x \log_2{x}=0,
\]
we can define the classical information entropy as
\[
S(x) = -p(x) \log_2{p(x)}.
\]

If we now consider a series of stochastic variables
$X=\{x_0,x_1,\dots,x_{n-1}\}$ with probability for an outcome $x\in X$
given by $p_X(x)$, the entropy becomes

\[
S_X = -\sum_{x\in X}p_X(x_i) \log_2{p_X(x_i)}.
\]
\end{frame}

\begin{frame}[plain,fragile]
\frametitle{Binary example}

If we now use the above binary example, we have (using that $p_1=p$ and $p_2=1-p$)
\[
S = -\sum_{i=1}^{2}p_i \log_2{p_i}=-p\log_2{p}-(1-p)\log_2{1-p}.
\]
For $p=0$ we have $S=0$ and for $p=1$ we get $S=0$. Thus, for the two
uniquely defined states we have entropy zero and a precise knowledge of the state. The maximum value is at
$p=0.5$ and this is a situation where our level of knowledge is no
better than us tossing a coin. This simple example demonstrates how we
can use the above mathematical expression as a way to say whether we
have a well-defined outcome or not. Our lack of knowledge is expressed
by the largest entropy value.
\end{frame}

\begin{frame}[plain,fragile]
\frametitle{Von Neumann entropy}

Using the density matrix $\rho$, we define the quantum mechanical equivalent of the classical entropy as 
\[
S=-\mathrm{Tr}[\rho\log_2{\rho}].
\]
This is the so-called Von Neumann entropy. How did we arrive at this expression? 

And why can't we use the classical definition for a quantum mechanical system as well?
\end{frame}

\begin{frame}[plain,fragile]
\frametitle{The density matrix again}

Let us assume we are studying specific system $A$. The density matrix,
using an ONB basis $\psi_j$, is defined as
\[
\bm{\rho}_A=\sum_j\lambda_j\vert \psi_j\rangle\langle \psi_j\vert=\sum_j p_j\vert \psi_j\rangle\langle \psi_j\vert,
\]
and the eigenvalues $\lambda_j$ are the probabilities (or overlap coefficients) of being in a
specific state $\psi_j$. 

Note that the density matrix/operator is a semi-positive definite matrix.
\end{frame}

\begin{frame}[plain,fragile]
\frametitle{Linking with a new expression for the entropy}

Using a unitary transformation $\bm{U}$ we can transform the density
matrix into a diagonal matrix $\bm{D}_a$ where the eigenvalues are the
above mentioned probabilities (overlap coefficients squared). We can then define
\[
\bm{D}_A=\bm{U}^{\dagger}\bm{\rho}_A\bm{U}=\begin{bmatrix} p_0 & 0 & 0 & \dots & 0\\
 0 & p_1 & 0 & \dots & 0\\
 0 & 0 & p_2 & \dots & 0\\
\dots & \dots & \dots & \dots & \dots\\
\dots & \dots & \dots & \dots & \dots\\
\dots & \dots & \dots & \dots & \dots\\
0 & 0 & \dots & 0 & p_{n-1}\end{bmatrix}.
\]
\end{frame}

\begin{frame}[plain,fragile]
\frametitle{Von Neumann entropy}

If we now construct the following object
\[
\bm{\rho}_A\log_2{\bm{\rho}_A},
\]
and perform a unitary transformation, recalling that $\bm{U}\bm{U}^{\dagger}=\bm{U}^{\dagger}\bm{U}=\bm{I}$,
\[
\bm{U}^{\dagger}\bm{\rho}_A\log_2{\bm{\rho}_A}\bm{U}=\bm{U}^{\dagger}\bm{\rho}_A\bm{I}\log_2{\bm{\rho}_A}\bm{U}=\bm{D}_A\bm{U}^{\dagger}\log_2{\bm{\rho}_A}\bm{U},
\]
and defining
\[
\bm{U}^{\dagger}\log_2{\bm{\rho}_A}\bm{U}=\log_2{\bm{D}_A},
\]
we have that 
\[
\bm{U}^{\dagger}\bm{\rho}_A\log_2{\bm{\rho}_A}\bm{U}=\bm{D}_A \log_2{\bm{D}_A}.
\]
\end{frame}

\begin{frame}[plain,fragile]
\frametitle{$\log_2{\bm{D}_A}$}

We have that
\[
\log_2{\bm{D}_A}=\begin{bmatrix} \log_2{p_0} & 0 & 0 & \dots & 0\\
 0 & \log_2{p_1} & 0 & \dots & 0\\
 0 & 0 & \log_2{p_2} & \dots & 0\\
\dots & \dots & \dots & \dots & \dots\\
\dots & \dots & \dots & \dots & \dots\\
\dots & \dots & \dots & \dots & \dots\\
0 & 0 & \dots & 0 & \log_2{p_{n-1}}\end{bmatrix}.
\]
\end{frame}

\begin{frame}[plain,fragile]
\frametitle{Entropy and traces}

If we take the trace of 
\[
\mathrm{Tr}[\bm{U}^{\dagger}\bm{\rho}_A\log_2{\bm{\rho}_A}\bm{U}]=\mathrm{Tr}[\bm{D}_A \log_2{\bm{D}_A}],
\]

which is nothing but
\[
\mathrm{Tr}[\bm{D}_A \log_2{\bm{D}_A}]=\sum_i p_i\log_2{p_i},
\]
which allows us to define the general Von-Neumann entropy as
\[
S_A=-\mathrm{Tr}[\bm{\rho}_A\log_2{\bm{\rho}_A}].
\]
Performing a unitary transformation gives the expression we know from classical information theory.
\end{frame}

\begin{frame}[plain,fragile]
\frametitle{What next?}

Before we proceed we are now going to study the entropy of a system
relevant for project 1.  We will compute the entropy of the ground
state and the entropy of a subsystem by tracing out (as done
previosly) the density matrix of ground state.

The system is a two-qubit and we will use this later to rewrite the
Hamiltonian matrix in terms of various Pauli-matrices.
\end{frame}

\begin{frame}[plain,fragile]
\frametitle{Two-qubit system and calculation of density matrices and project 1}

\textbf{This part is best seen using the jupyter-notebook}.

This system can be thought of as composed of two subsystems
$A$ and $B$. Each subsystem has computational basis states

\[
\vert 0\rangle_{\mathrm{A,B}}=\begin{bmatrix} 1 & 0\end{bmatrix}^T \hspace{1cm} \vert 1\rangle_{\mathrm{A,B}}=\begin{bmatrix} 0 & 1\end{bmatrix}^T.
\]
The subsystems could represent single particles or composite many-particle systems of a given symmetry.
This leads to the many-body computational basis states

\[
\vert 00\rangle = \vert 0\rangle_{\mathrm{A}}\otimes \vert 0\rangle_{\mathrm{B}}=\begin{bmatrix} 1 & 0 & 0 &0\end{bmatrix}^T,
\]
and
\[
\vert 01\rangle = \vert 0\rangle_{\mathrm{A}}\otimes \vert 1\rangle_{\mathrm{B}}=\begin{bmatrix} 0 & 1 & 0 &0\end{bmatrix}^T,
\]
and
\[
\vert 10\rangle = \vert 1\rangle_{\mathrm{A}}\otimes \vert 0\rangle_{\mathrm{B}}=\begin{bmatrix} 0 & 0 & 1 &0\end{bmatrix}^T,
\]
and finally
\[
\vert 11\rangle = \vert 1\rangle_{\mathrm{A}}\otimes \vert 1\rangle_{\mathrm{B}}=\begin{bmatrix} 0 & 0 & 0 &1\end{bmatrix}^T.
\]

These computational basis states define the eigenstates of the non-interacting  Hamiltonian
\[
H_0\vert 00 \rangle = \epsilon_{00}\vert 00 \rangle,
\]
\[
H_0\vert 10 \rangle = \epsilon_{10}\vert 10 \rangle,
\]
\[
H_0\vert 01 \rangle = \epsilon_{01}\vert 01 \rangle,
\]
and
\[
H_0\vert 11 \rangle = \epsilon_{11}\vert 11 \rangle.
\]
The interacting part of the Hamiltonian $H_{\mathrm{I}}$ is given by the tensor product of two $\sigma_x$ and $\sigma_z$  matrices, respectively, that is
\[
H_{\mathrm{I}}=H_x\sigma_x\otimes\sigma_x+H_z\sigma_z\otimes\sigma_z,
\]
where $H_x$ and $H_z$ are interaction strength parameters. Our final Hamiltonian matrix is given by
\[
\bm{H}=\begin{bmatrix} \epsilon_{00}+H_z & 0 & 0 & H_x \\
                       0  & \epsilon_{10}-H_z & H_x & 0 \\
		       0 & H_x & \epsilon_{01}-H_z & 0 \\
		       H_x & 0 & 0 & \epsilon_{11} +H_z \end{bmatrix}.
\] 

The four eigenstates of the above Hamiltonian matrix can in turn be used to
define density matrices. As an example, the density matrix of the
first eigenstate (lowest energy $E_0$) $\Psi_0$ is given by the outerproduct

\[
\rho_0=\left(\alpha_{00}\vert 00 \rangle+\alpha_{10}\vert 10 \rangle+\alpha_{01}\vert 01 \rangle+\alpha_{11}\vert 11 \rangle\right)\left(\alpha_{00}^*\langle 00\vert+\alpha_{10}^*\langle 10\vert+\alpha_{01}^*\langle 01\vert+\alpha_{11}^*\langle 11\vert\right),
\]

where the coefficients $\alpha_{ij}$ are the eigenvector coefficients
resulting from the solution of the above eigenvalue problem. 

We can
then in turn define the density matrix for the subsets $A$ or $B$ as

\[
\rho_A=\mathrm{Tr}_B(\rho_{0})=\langle 0 \vert \rho_{0} \vert 0\rangle_{B}+\langle 1 \vert \rho_{0} \vert 1\rangle_{B},
\]

or

\[
\rho_B=\mathrm{Tr}_A(\rho_0)=\langle 0 \vert \rho_{0} \vert 0\rangle_{A}+\langle 1 \vert \rho_{0} \vert 1\rangle_{A}.
\]

The density matrices for these subsets can be used to compute the
so-called von Neumann entropy, which is one of the possible measures
of entanglement. A pure state has entropy equal zero while entangled
state have an entropy larger than zero. The von-Neumann entropy is
defined as

\[
S(A,B)=-\mathrm{Tr}\left(\rho_{A,B}\log_2 (\rho_{A,B})\right).
\]

The example here shows the above von Neumann entropy based on the
density matrix for the lowest many-body state. We see clearly a jump
in the entropy around the point where we have a level crossing. At
interaction strenght $\lambda=0$ we have many-body states purely
defined by their computational basis states. As we switch on the
interaction strength, we obtain an increased degree of mixing and the
entropy increases till we reach the level crossing point where we see
an additional and sudden increase in entropy. Similar behaviors are
observed for the other states. The most important result from this
example is that entanglement is driven by the Hamiltonian itself and
the strength of the interaction matrix elements and the
non-interacting energies.
































































\begin{minted}[fontsize=\fontsize{9pt}{9pt},linenos=false,mathescape,baselinestretch=1.0,fontfamily=tt,xleftmargin=2mm]{python}
%matplotlib inline
from  matplotlib import pyplot as plt
import numpy as np
from scipy.linalg import logm, expm
def log2M(a): # base 2 matrix logarithm
    return logm(a)/np.log(2.0)

dim = 4
Hamiltonian = np.zeros((dim,dim))
#number of lambda values
n = 40
lmbd = np.linspace(0.0,1.0,n)
Hx = 2.0
Hz = 3.0
# Non-diagonal part as sigma_x tensor product with sigma_x
sx = np.matrix([[0,1],[1,0]])
sx2 = Hx*np.kron(sx, sx)
# Diagonal part as sigma_z tensor product with sigma_z
sz = np.matrix([[1,0],[0,-1]])
sz2 = Hz*np.kron(sz, sz)
noninteracting = [0.0, 2.5, 6.5, 7.0]
D = np.diag(noninteracting)
Eigenvalue = np.zeros((dim,n))
Entropy = np.zeros(n)

for i in range(n): 
    Hamiltonian = lmbd[i]*(sx2+sz2)+D
    # diagonalize and obtain eigenvalues, not necessarily sorted
    EigValues, EigVectors = np.linalg.eig(Hamiltonian)
    # sort eigenvectors and eigenvalues
    permute = EigValues.argsort()
    EigValues = EigValues[permute]
    EigVectors = EigVectors[:,permute]
    # Compute density matrix for selected system state, here ground state
    DensityMatrix = np.zeros((dim,dim))
    DensityMatrix = np.outer(EigVectors[:,0],EigVectors[:,0])
    # Project down on substates and find density matrix for subsystem
    d = np.matrix([[1,0],[0,1]])
    v1 = [1.0,0.0]
    proj1 = np.kron(v1,d)
    x1 = proj1 @ DensityMatrix @ proj1.T
    v2 = [0.0,1.0]
    proj2 = np.kron(v2,d)
    x2 = proj2 @ DensityMatrix @ proj2.T
    # Total density matrix for subsystem
    total = x1+x2
    # von Neumann Entropy for subsystem 
    Entropy[i] = -np.matrix.trace(total @ log2M(total))
    # Plotting eigenvalues and entropy as functions of interaction strengths
    Eigenvalue[0,i] = EigValues[0]
    Eigenvalue[1,i] = EigValues[1]
    Eigenvalue[2,i] = EigValues[2]
    Eigenvalue[3,i] = EigValues[3]
plt.plot(lmbd, Eigenvalue[0,:] ,'b-',lmbd, Eigenvalue[1,:],'g-',)
plt.plot(lmbd, Eigenvalue[2,:] ,'r-',lmbd, Eigenvalue[3,:],'y-',)
plt.xlabel('$\lambda$')
plt.ylabel('Eigenvalues')
plt.show()
plt.plot(lmbd, Entropy)
plt.xlabel('$\lambda$')
plt.ylabel('Entropy')          
plt.show

\end{minted}
\end{frame}

\begin{frame}[plain,fragile]
\frametitle{Schmidt decomposition}

A way to study entanglement is through the so-called Schmidt decomposition,
which is essentially an application of the
singular-value decomposition.
\end{frame}

\begin{frame}[plain,fragile]
\frametitle{Pure states and Schmidt decomposition}

The Schmidt decomposition allows us to define a pure state in a
bipartite Hilbert space composed of systems $A$ and $B$ as

\[
\vert\psi\rangle=\sum_{i=0}^{d-1}\sigma_i\vert i\rangle_A\vert i\rangle_B,
\]
where the amplitudes $\sigma_i$ are real and positive and their
squared values sum up to one, $\sum_i\sigma_i^2=1$. The states $\vert
i\rangle_A$ and $\vert i\rangle_B$ form orthornormal bases for systems
$A$ and $B$ respectively, the amplitudes $\lambda_i$ are the so-called
Schmidt coefficients and the Schmidt rank $d$ is equal to the number
of Schmidt coefficients and is smaller or equal to the minimum
dimensionality of system $A$ and system $B$, that is $d\leq
\mathrm{min}(\mathrm{dim}(A), \mathrm{dim}(B))$.
\end{frame}

\begin{frame}[plain,fragile]
\frametitle{Proof of Schmidt decomposition}

The proof for the above decomposition is based on the singular-value
decomposition. To see this, assume that we have two orthonormal bases
sets for systems $A$ and $B$, respectively. That is we have two ONBs
$\vert i\rangle_A$ and $\vert j\rangle_B$. We can always construct a
product state (a pure state) as

\[ \vert\psi \rangle = \sum_{ij} c_{ij}\vert i\rangle_A\vert
j\rangle_B,
\]
where the coefficients $c_{ij}$ are the overlap coefficients which
belong to a matrix $\bm{C}$.
\end{frame}

\begin{frame}[plain,fragile]
\frametitle{Further parts of proof}

If we now assume that the
dimensionalities of the two subsystems $A$ and $B$ are the same $d$,
we can always rewrite the matrix $\bm{C}$ in terms of a singular-value
decomposition with unitary/orthogonal matrices $\bm{U}$ and $\bm{V}$
of dimension $d\times d$ and a matrix $\bm{\Sigma}$ which contains the
(diagonal) singular values $\sigma_0\leq \sigma_1 \leq \dots 0$ as

\[
\bm{C}=\bm{U}\bm{\Sigma}\bm{V}^{\dagger}.
\]
\end{frame}

\begin{frame}[plain,fragile]
\frametitle{SVD parts in proof}

This means we can rewrite the coefficients $c_{ij}$ in terms of the singular-value decomposition
\[
c_{ij}=\sum_k u_{ik}\sigma_kv_{kj},
\]
and inserting this in the definition of the pure state $\vert \psi\rangle$ we have

\[
\vert\psi \rangle = \sum_{ij} \left(\sum_k u_{ik}\sigma_kv_{kj} \right)\vert i\rangle_A\vert j\rangle_B.
\]
\end{frame}

\begin{frame}[plain,fragile]
\frametitle{Slight rewrite}

We rewrite the last equation as

\[
\vert\psi \rangle = \sum_{k}\sigma_k \left(\sum_i u_{ik}\vert i\rangle_A\right)\otimes\left(\sum_jv_{kj}\vert j\rangle_B\right),
\]
which we identify simply as, since the matrices $\bm{U}$ and $\bm{V}$ represent unitary transformations,
\[
\vert\psi \rangle = \sum_{k}\sigma_k \vert k\rangle_A\vert k\rangle_B.
\]
\end{frame}

\begin{frame}[plain,fragile]
\frametitle{Different dimensionalities}

It is straight forward to prove this relation in case systems $A$ and
$B$ have different dimensionalities.  Once we know the Schmidt
decomposition of a state, we can immmediately say whether it is
entangled or not. If a state $\psi$ has is entangled, then its Schmidt
decomposition has more than one term. Stated differently, the state is
entangled if the so-called Schmidt rank is is greater than one.  There
is another important property of the Schmidt decomposition which is
related to the properties of the density matrices and their trace
operations and the entropies. In order to introduce these concepts let
us look at the two-qubit Hamiltonian described here.
\end{frame}

\begin{frame}[plain,fragile]
\frametitle{Rewriting the above Hamiltonian as a set of Pauli-gates}

Last week we discussed how to rewrite a one-qubit Hamiltonian in terms
of various one-qubit gates.  This forms our first Hamiltonian of
project 1. We will move on to a two-qubit system.  But first a short
reminder about specific two-qubit gates.
\end{frame}

\begin{frame}[plain,fragile]
\frametitle{Two-Qubit Gates}

A two-qubit gate is a physical action that is applied to two
qubits. It can be represented by a matrix $U$ from the group
SU(4). One important type of two-qubit gates are controlled gates,
which work as follows: Suppose $U$ is a single-qubit gate. A
controlled-$U$ gate ($CU$) acts on two qubits: a control qubit
$\vert x \rangle $ and a target qubit $\vert y \rangle $. The controlled-$U$ gate
applies the identity $I$ or the single-qubit gate $U$ to the target
qubit if the control gate is in the zero state $\vert 0\rangle$ or the one
state $\vert 1\rangle$, respectively.
\end{frame}

\begin{frame}[plain,fragile]
\frametitle{Control qubit}

The control qubit is not acted
upon. This can be represented as follows if
\[CU\vert xy\rangle=
\vert xy\rangle \hspace{0.1cm} \mathrm{if} \hspace{0.1cm}  \vert x \rangle =\vert 0\rangle.
\]
\end{frame}

\begin{frame}[plain,fragile]
\frametitle{In matrix form}

It is easier to see in a matrix form.
It can be written in matrix form by writing it as a superposition of
the two possible cases, each written as a simple tensor product

\[
CU = \vert 0\rangle\langle 0\vert\otimes I + \vert 1\rangle\langle 1 \vert \otimes U=\begin{bmatrix}
1 & 0 & 0 & 0 \\
0 & 1 & 0 & 0 \\
0 & 0 & u_{00} & u_{01} \\
0 & 0 & u_{10} & u_{11}
\end{bmatrix}.
\]
\end{frame}

\begin{frame}[plain,fragile]
\frametitle{CNOT gate}

One of the most fundamental controlled gates is the CNOT gate. It is
defined as the controlled-$X$ gate $CX$. It can be written in matrix form as follows:

\[
\mathrm{CNOT}=\mathrm{CX}=\begin{bmatrix}
1 & 0 & 0 & 0 \\
0 & 1 & 0 & 0 \\
0 & 0 & 0 & 1 \\
0 & 0 & 1 & 0
\end{bmatrix}.
\]
\end{frame}

\begin{frame}[plain,fragile]
\frametitle{$\mathrm{CX}$ gate}

It changes, when operating on a two-qubit state where the first qubit is the control qubit and the second qubit is the target qubit, the states (check this)
\[
\mathrm{CX}\vert 00\rangle=\vert 00\rangle,
\]
\[
\mathrm{CX}\vert 10\rangle= \vert 11\rangle,
\]
\[
\mathrm{CX}\vert 01\rangle= \vert 01\rangle,
\]
\[
\mathrm{CX}\vert 11\rangle= \vert 10\rangle,
\]
which you can easily see by simply multiplying the above matrix with any of the above states.
\end{frame}

\begin{frame}[plain,fragile]
\frametitle{Swap gate}

A widely used two-qubit gate that goes beyond the simple controlled function is the SWAP gate. It swaps the states of the two qubits it acts upon

\[
\mathrm{SWAP}\vert xy\rangle=\vert yx\rangle.
\]
and has the following matrix form

\[
\mathrm{SWAP}
=\begin{bmatrix}
1 & 0 & 0 & 0 \\
0 & 0 & 1 & 0 \\
0 & 1 & 0 & 0 \\
0 & 0 & 0 & 1
\end{bmatrix}.
\]
\end{frame}

\begin{frame}[plain,fragile]
\frametitle{Reminder from last week}

We rewrote the one-qubit Hamiltonian
\[
\bm{H}=\begin{bmatrix} \epsilon_{1}+V_{11} & V_{12}\\
                       V_{21}  & \epsilon_{2}-V_{22} \end{bmatrix},
\] 
as (defining $V_{12}=V_{21}=V_x$) 
\[
\bm{H}=\frac{\epsilon_{1}+\epsilon_{2}}{2}\bm{I}+\frac{\epsilon_{1}-\epsilon_{2}}{2}\bm{Z}+\frac{V_{11}+V_{22}}{2}\bm{I}+\frac{V_{11}-V_{22}}{2}\bm{Z}+V_x\bm{X}.
\] 
We will now rewrite our first two-qubit Hamiltonian along the same lines.
\end{frame}

\begin{frame}[plain,fragile]
\frametitle{Two-qubit Hamiltonian}

We start again with 
\[
\bm{H}=\begin{bmatrix} \epsilon_{1}+V_z & 0 & 0 & V_x \\
                       0  & \epsilon_{2}-V_z & V_x & 0 \\
		       0 & H_x & \epsilon_{3}-V_z & 0 \\
		       H_x & 0 & 0 & \epsilon_{4} +V_z \end{bmatrix}.
\] 

This Hamiltonian, as discussed in the first exercise can be rewritten in terms of various one-qubit matrices.
\end{frame}

\begin{frame}[plain,fragile]
\frametitle{Definitions}

We define
\[
\epsilon_{II}=\frac{\epsilon_{1}+\epsilon_{2}+\epsilon_{3}+\epsilon_{4}}{4},
\]
\[
\epsilon_{ZI}=\frac{\epsilon_{1}+\epsilon_{2}-\epsilon_{3}-\epsilon_{4}}{4},
\]
\[
\epsilon_{IZ}=\frac{\epsilon_{1}-\epsilon_{2}+\epsilon_{3}-\epsilon_{4}}{4},
\]
\[
\epsilon_{ZZ}=\frac{\epsilon_{1}-\epsilon_{2}-\epsilon_{3}+\epsilon_{4}}{4}.
\]
\end{frame}

\begin{frame}[plain,fragile]
\frametitle{Ex1: Two-qubit Hamiltonian}

This exercise is also part of the first project.
Rewrite the Hamiltonian in term of various Pauli matrices and show that you can rwrite the Hamiltonian as
\[
\bm{H}=\bm{H}_0+\bm{H}_I,
\]
with 
\[
\bm{H}_0=\epsilon_{II}\bm{I}\otimes\bm{I}+\epsilon_{ZI}\bm{Z}\otimes\bm{I}+\epsilon_{IZ}\bm{I}\otimes\bm{Z}+\epsilon_{ZZ}\bm{Z}\otimes\bm{Z}
\]
and 
\[
\bm{H}_I=V_x\bm{X}\otimes\bm{X}+V_z\bm{Z}\otimes\bm{Z}
\]
\end{frame}

\begin{frame}[plain,fragile]
\frametitle{Ex2: Two-qubit Hamiltonian and entropy}

Use the Hamiltonian for the two-qubit example to find the eigenpairs
as functions of the interaction strength $\lambda$ and study the final
eigenvectors as functions of the admixture of the original basis
states.  Discuss the results as functions of the parameter $\lambda$ and compute the von Neumann
entropy and discuss the results. You will need to calculate the entropy of the subsystems $A$ or $B$.
\end{frame}

\end{document}
